% Source: https://artofproblemsolving.com/community/c3h1602150_prime_numbers
% Stable identity: aops-topic-1602150-q01
Write $S(p)=p+p^2+\cdots+p^{10}$, so $q=S(p)-2017$. The primes below $10$ are $2,3,5,7$.

\textbf{$p=2$:} $S(2)=2^{11}-2=2046$, so $q=29$, which is prime.

\textbf{$p=7$:} since $7\equiv1\pmod3$, every power $7^k\equiv1\pmod 3$, so $S(7)\equiv10\equiv1$, and $2017\equiv1\pmod3$. Hence $3\mid q$; but $q=S(7)-2017$ is much larger than $3$, so $q$ is not prime.

\textbf{$p=3$:} $S(3)=\dfrac{3^{11}-3}{2}=88572$, so $q=86555$, which ends in $5$ and is not prime.

\textbf{$p=5$:} $S(5)=\dfrac{5^{11}-5}{4}=12207030$, so $q=12205013=223\cdot229\cdot239$, which is not prime.

So the only pair is $\boxed{(p,q)=(2,29)}$.

\emph{Remark.} Without the condition $p<10$ the answer is not unique: the same argument shows that $p\equiv1\pmod3$ never works, but some larger primes $p\equiv2\pmod 3$ do give a prime $q$; for example $p=197$ gives the prime $q=88485562579618578299669$ (checked by computer).
